\chapter{Workflow di ottimizzazione}
\label{ch:workflow_di_ottimizzazione}

\grokfig{fig_time_budget_clessidra}{La clessidra del time budget:
sabbia a strati colorati ripartisce il tempo fra Phase~A, Phase~B
e metaeuristiche.}

\section{Schema delle fasi}

\begin{figure}[h]
\centering
\begin{tikzpicture}[node distance=0.5cm and 0.6cm,
  rect/.style={draw, rounded corners=4pt, very thick, fill=cBox!30,
               minimum width=3.4cm, minimum height=1.2cm,
               align=center, font=\small}]
  \node[rect] (a) {Phase A\\\textit{assignment}\\CP-SAT};
  \node[rect, right=of a] (b) {Phase B\\\textit{scheduling}\\+ decomp. spettrale};
  \node[rect, right=of b] (m) {Metaeuristiche\\\textit{LNS+SA+TS+ILS}};
  \node[rect, right=of m] (r) {Phase 4\\\textit{aule}\\CP-SAT};
  \draw[-{Stealth[length=3mm]}, thick] (a) -- (b);
  \draw[-{Stealth[length=3mm]}, thick] (b) -- (m);
  \draw[-{Stealth[length=3mm]}, thick] (m) -- (r);

  \node[below=0.4cm of a, font=\scriptsize] {assignments};
  \node[below=0.4cm of b, font=\scriptsize] {lessons};
  \node[below=0.4cm of m, font=\scriptsize] {lessons SOFT-ottime};
  \node[below=0.4cm of r, font=\scriptsize] {lessons.classroom};
\end{tikzpicture}
\caption{Pipeline a 4 fasi.}
\label{fig:pipeline}
\end{figure}

\section{Phase A: assegnazione (cattedre)}

\paragraph{Modulo.} \texttt{experiments/cpsat\_v2\_assignment.py}.
\paragraph{Lancio.} \texttt{POST /api/optimize/assignment} con
\texttt{\{time\_limit\_s, workers, log\}}.

\paragraph{Input.} L'insieme di \texttt{(class, subject,
hours\_per\_week)} letto da \texttt{class\_subjects} e/o
\texttt{curriculum\_subject\_hours} + le abilitazioni docente in
\texttt{teacher\_subjects} + la classe-di-concorso preference in
\texttt{subject\_group\_weights}.

\paragraph{Output.} La tabella \texttt{assignments}
(\texttt{teacher x class x subject} con \texttt{hours}). Problema di
covering: ogni \texttt{(class, subject)} riceve un docente abilitato,
rispettando max-hours per docente, classi-di-concorso e
\texttt{teacher\_compatible\_classes}.

\section{Phase B: scheduling}

\paragraph{Modulo.} \texttt{experiments/cpsat\_v2\_timetable.py} +
\texttt{experiments/decomposition\_spectral\_v2.py}.

\paragraph{Decomposizione spettrale.}
Per istanze grandi, il problema CP-SAT integrale \`e intrattabile.
La pipeline calcola un grafo di co-occupazione fra docenti, ne fa
un'embedding spettrale via Laplaciano normalizzato e usa K-means su
quell'embedding per partizionare in $k$ cluster
\emph{loosely-coupled}. Ogni cluster diventa una sub-istanza CP-SAT
indipendente; un secondo step "bridges" risolve solo i bordi; un
terzo step "ricucitura" fa full-CPSAT con tutto il problema partendo
dalla soluzione stitched.

\paragraph{Scope del solver CP-SAT (\texttt{cp\_sat\_scope}).}
Nella card 3 di \texttt{/optimize} il dropdown \emph{Scope del
solver CP-SAT} sceglie come la Phase B passa il problema al
risolutore:
\begin{itemize}
  \item \textbf{Per giorno} (\texttt{day}, default).
        Il solver risolve un giorno alla volta. La distribuzione
        settimanale delle ore di ciascuna cattedra (\emph{day count})
        viene precalcolata da Phase A e iniettata come uguaglianza
        rigida per ogni giorno. Veloce, scalabile fino a scuole grandi.
  \item \textbf{Settimana intera} (\texttt{week}).
        Il solver \texttt{MonolithicSolver(scope=None)} vede tutta
        la settimana in un solo modello CP-SAT. Pi\`u robusto sui
        vincoli che attraversano i giorni (compresenze, sostegno,
        gruppi a geometria variabile) ma significativamente pi\`u
        lento: i workers devono esplorare uno spazio di ricerca di
        sei volte la cardinalit\`a del per-giorno.
\end{itemize}

\paragraph{Modalit\`a Phase A (\texttt{phase\_a\_mode}).}
Phase A \`e il pre-step che decide quante ore di ciascuna cattedra
cadono su ciascun giorno della settimana (\emph{day count}). Il
secondo dropdown della card 3 controlla come questo output viene
combinato con il solver:
\begin{itemize}
  \item \textbf{Sempre} (\texttt{always}, default per scope = Per
        giorno). Phase A esegue e il \texttt{day\_count} viene
        imposto come uguaglianza rigida. Obbligatorio in modalit\`a
        Per giorno -- senza \texttt{day\_count} il solver per-giorno
        non sa quante ore sono "previste" in quel giorno.
  \item \textbf{Salta} (\texttt{skip}, solo con scope = Settimana
        intera). Phase A non viene eseguita; \`e la Settimana intera
        a decidere autonomamente come distribuire le ore. Massima
        libert\`a per il solver, costo: nessun warm start.
  \item \textbf{Hint morbido} (\texttt{soft\_hint}, solo con scope
        = Settimana intera). Phase A esegue come al solito ma il
        suo output viene iniettato nel solver tramite
        \texttt{model.AddHint} sui contatori per-giorno -- il solver
        parte dalla distribuzione di Phase A come prima soluzione
        ammissibile e pu\`o deviare se trova qualcosa di meglio.
        Tipicamente il compromesso preferibile per scuole complesse:
        si tiene il warm start, si toglie la rigidit\`a.
\end{itemize}

Vincolo incrociato (validato sia dal frontend sia dal backend):
\texttt{day} richiede \texttt{always}; \texttt{week} esclude
\texttt{always}. La UI disabilita le combinazioni invalide e ricorregge
automaticamente \texttt{phase\_a\_mode} quando l'utente cambia
\texttt{cp\_sat\_scope}.

\paragraph{Quando scegliere cosa.}
\begin{itemize}
  \item \emph{Scuola standard, focus sulla velocit\`a}: \texttt{day
        + always}. \`E il default per qualcosa: scala bene e produce
        soluzioni di buona qualit\`a in tempo lineare nel numero di
        classi.
  \item \emph{Scuola piccola con molte compresenze e sostegno}:
        \texttt{week + skip}. La rigidit\`a del \texttt{day\_count}
        di Phase A pu\`o rendere infeasibile un orario che invece
        esiste; lasciar decidere alla Settimana intera evita il
        problema.
  \item \emph{Scuola grande ma con vincoli incrociati notevoli}:
        \texttt{week + soft\_hint}. Il warm start velocizza la
        ricerca e la Settimana intera resta libera di adattare la
        distribuzione.
\end{itemize}

\paragraph{Branch-and-Price V1 e \texttt{cp\_sat\_scope=week}.}
Una domanda emersa durante lo sviluppo della Phase 3 era: se
Phase A non fissa pi\`u il \texttt{day\_count} come uguaglianza
rigida (\texttt{phase\_a\_mode=skip} oppure \texttt{soft\_hint}),
la generazione di colonne V1 -- master con vincoli di copertura
indicizzati su \texttt{(teacher, class, subject, day)} -- riesce a
iterare di pi\`u della singola passata
\texttt{no\_improving\_column} osservata in modalit\`a
\texttt{day}? L'intuizione era che un \texttt{dc\_value} ``allentato''
generasse duali $\lambda$ non-zero, sbloccando colonne migliorative.

La risposta empirica, misurata da
\texttt{tests/benchmarks/run\_bp\_v1\_scope\_week.py} su due scenari
sintetici (\texttt{small} = 10 docenti / 3 classi, \texttt{medium}
= 4 docenti / 3 classi con cattedre dense), \`e \emph{no}: le
cinque granularit\`a V1 (\texttt{teacher}, \texttt{teacher-day},
\texttt{teacher-class}, \texttt{teacher-class-subject},
\texttt{teacher-subject}) terminano sempre alla prima iterazione
con \texttt{bp\_terminated\_reason=no\_improving\_column} e
\texttt{bp\_columns\_added\_total=0}, sia quando il
\texttt{dc\_value} viene da Phase A sia quando viene derivato dalla
soluzione del solver settimanale.

\begin{itemize}
  \item \texttt{small}: \texttt{dc\_phase\_a} (21 chiavi non-zero)
        differisce strutturalmente da \texttt{dc\_week} (16 chiavi
        non-zero) -- 17 chiavi solo in $A$, 12 solo in $B$, 2
        chiavi in comune con valori diversi -- ma per ogni
        granularit\`a entrambe le sorgenti producono
        \texttt{iters=1, cols=0}. La soluzione \`e HARD-feasible
        in entrambi i casi; il soft score finale \`e $240$ con
        \texttt{phase\_a} e $0$ con \texttt{week}, ma la
        differenza viene dal seed pattern, non dalle iterazioni
        BP.
  \item \texttt{medium}: \texttt{dc\_phase\_a} (69 chiavi) vs
        \texttt{dc\_week} (66 chiavi), 34 vs 31 chiavi disgiunte,
        10 chiavi in comune con valori diversi. Stesso esito:
        \texttt{iters=1, cols=0} su tutte le granularit\`a, sia
        con \texttt{phase\_a} sia con \texttt{week}.
\end{itemize}

\noindent
La spiegazione \`e strutturale, non un bug della Phase 3. I seed
pattern di \texttt{column\_generation.\_seed\_patterns} sono
greedy-ottimali rispetto al \texttt{dc\_value} fornito: per ogni
docente piazzano esattamente le sue ore nei giorni indicati, in
ordine \texttt{(class, subject, day, hour)} crescente, evitando
sovrapposizioni. Il master LP V1 sceglie il pattern con peso
maggiore -- tipicamente il primo seed -- e i duali $\lambda$ del
sub-problema di pricing risultano gi\`a saturati: il pricer
ricostruisce la stessa griglia di ore e non trova \emph{reduced
cost} negativa. Cambiare la sorgente del \texttt{dc\_value}
modifica la \emph{distribuzione} delle ore (giorni e quantit\`a)
ma non cambia il fatto che ciascun docente ha un'unica ``forma''
greedy-coerente con la propria domanda. Per sbloccare iterazioni
multiple servirebbe (a) generare seed pattern \emph{strutturalmente
diversi} (perturbazioni randomiche delle ore o degli ordinamenti),
(b) introdurre vincoli di accoppiamento cross-teacher gi\`a nel
master (variante V2 con cover indicizzato su \texttt{(class, day)}
o \texttt{(curriculum, day)}), oppure (c) iniettare una pricer
heuristic con duali \emph{stabilizzati} che esplori intorno
all'ottimo locale del seed. Sono tre direzioni di lavoro distinte
e fuori dallo scope di Phase 3-5.

In pratica: \texttt{cp\_sat\_scope=week} sblocca il
\emph{solver settimanale} (la distribuzione delle ore non \`e
pi\`u prerogativa di Phase A) ma NON sblocca il loop BP V1, che
resta una single-shot ottimizzazione del seed catalog. Per orario
di buona qualit\`a su scuole grandi con BP, la leva attuale resta
la granularit\`a V2 (\texttt{class}, \texttt{class-day},
\texttt{day}, \texttt{curriculum}) o l'\texttt{iterative-diversified}
classico, non l'interazione fra scope settimanale e V1.

\section{Metaeuristiche}

\paragraph{Modulo.} \texttt{experiments/metaheuristics.py}.
Cascata di:

\begin{itemize}
  \item \textbf{LNS} -- destroy-and-repair su finestre
        \texttt{(teacher, day)} o \texttt{(class, day)}; freeze del
        60-70\% delle lezioni e CP-SAT-resolve del resto.
  \item \textbf{SA} (Simulated Annealing) -- single-swap accettati
        con Metropolis a $T$ decrescente con fattore $\alpha$.
  \item \textbf{TS} (Tabu Search) -- best-improvement local search
        con tabu list (default size 80).
  \item \textbf{ILS} (Iterated Local Search) -- alterna $n$ cicli di
        TS con \emph{kick} perturbativi (random 4-opt fra giorni).
\end{itemize}

Tutti contribuiscono al SOFT score: holes per teacher, isolated 1-
and 5-hour days, weekly distribution, dual-hour pairs, no 6th hour,
banded preferences su materie. Il SOFT include la contribuzione delle
matrici di disponibilit\`a SOFT/PREFERRED e dei logical SOFT/PREFERRED.

\section{Phase 4: assegnazione aule}

\paragraph{Modulo.} \texttt{experiments/classroom\_assignment.py}.
Input: la tabella \texttt{lessons} (g\`ia temporalmente fissata) +
le classroom rules (\texttt{kind}, \texttt{capacity},
\texttt{multi\_class}, preferenze materia$\leftrightarrow$aula,
classe$\leftrightarrow$aula, docente$\leftrightarrow$aula).
Output: \texttt{lessons.classroom\_name} popolato.

\section{Drag-and-drop con preview live}

Trascinando una lezione su un altro slot la UI fa due chiamate:

\begin{enumerate}
  \item \texttt{POST /api/schedule/move-preview} -- il backend simula
        la mossa per tutte le 36 (day, hour) e per ognuna restituisce:
        \begin{itemize}
          \item \texttt{ok} con \texttt{delta\_soft <= 0} -- mossa
                accettabile (verde se delta $<$ 0, verde chiaro
                se delta $=$ 0).
          \item \texttt{soft\_worse} -- mossa accettabile ma peggiora
                il SOFT (giallo).
          \item \texttt{hard\_violation} -- vincolo HARD violato
                (rosso).
          \item \texttt{noop} -- slot di origine.
        \end{itemize}
  \item L'utente droppa, la UI invia \texttt{PUT
        /api/schedule/move-lesson} che valida e persiste.
\end{enumerate}

\subsection{Room-follows-lesson}

Il move-preview NON considera l'aula nel calcolo HARD. Se la nuova
posizione \`e compatibile con docente / classe ma l'aula attuale
\`e g\`ia occupata o HARD-unavailable in destination, il backend:

\begin{enumerate}
  \item Sposta la lezione.
  \item Lascia l'aula libera sul nuovo slot solo se non c'\`e conflitto.
  \item Se c'\`e conflitto, rimuove l'aula dalla lezione spostata e
        risponde con \texttt{room\_cleared=true, cleared\_room=<nome>}.
  \item La UI apre un Modal che spiega cosa \`e successo e chiede di
        scegliere una nuova aula dal dropdown (verde libera / rosso
        occupata).
\end{enumerate}

\section{Slot picker matriciale (Monitor)}

Nel tab \emph{Monitor}, espandendo una cattedra si vedono le singole
lezioni; cliccando il bottone Giorno/Ora si apre un modal con una
matrice 6x6 dove ogni cella \`e colorata in base alla disponibilit\`a:

\begin{itemize}
  \item verde: free (docente e classe liberi)
  \item rosso: HARD (docente o classe g\`ia impegnati)
  \item ambra: aula occupata da un'altra lezione
  \item azzurro: slot attuale
\end{itemize}

Click su una cella libera fa un dry-run via
\texttt{PUT /api/monitor/event/\{aid\}/lesson/\{lid\}}. Se
ci sono conflitti, apre un secondo Modal con 3 opzioni:
\textbf{Annulla}, \textbf{Disassegna le lezioni in conflitto},
\textbf{Disassegna e ottimizza dopo}.

\section{Punteggio di graduatoria e criterio Anzianita}

Ciascun docente nell'anagrafica ha un campo
\textbf{graduatoria\_score} (numerico). Il valore esprime il
punteggio di graduatoria utile a discriminare i docenti quando
si tratta di assegnare cattedre o supplenze: pi\`u alto =
maggiore precedenza.

Il campo \`e usato in due punti:

\begin{enumerate}
\item \textbf{Modulo supplenze}: quando il sistema cerca un
  candidato per coprire un'assenza, ordina la rosa per
  \texttt{graduatoria\_score} decrescente.
\item \textbf{Phase A, criterio ``Anzianita''}: il DSL della
  funzione obiettivo della Phase A include un preset
  \textbf{Anzianita} che, oltre ai criteri standard di bilancio
  (capacit\`a non utilizzata, n. di indirizzi, ecc.),
  privilegia le assegnazioni che fanno coincidere docenti con
  alto \texttt{graduatoria\_score} con classi del triennio
  finale (a parit\`a di altre condizioni).
\end{enumerate}

Il preset \`e selezionabile dalla card \textbf{Phase A
(assegnazione)} del tab Workflow, dropdown ``Criterio''. Il DSL
completo della funzione obiettivo \`e visibile nel modal
``Mostra DSL attivo'' della stessa card.

% =========== Cap 6: UI guide ===========
